\documentclass[twocolumn]{article}

% Language setting
% Replace `english' with e.g. `spanish' to change the document language
\usepackage[spanish,es-lcroman]{babel}

% Set page size and margins
% Replace `letterpaper' with `a4paper' for UK/EU standard size
\usepackage[letterpaper,top=2cm,bottom=2cm,left=2cm,right=2cm,marginparwidth=1.75cm]{geometry}

% Useful packages
\usepackage{amsthm}
\usepackage{amssymb}
\usepackage{amsmath}
\usepackage{amsfonts}
\usepackage{mathtools}
\usepackage{enumitem}
\usepackage{graphicx}
\usepackage{subcaption}
\usepackage[colorlinks=true, allcolors=blue]{hyperref}
\usepackage[utf8]{inputenc}
\usepackage{lipsum} % para generar texto de prueba
\usepackage{xcolor} % Para el color
\usepackage{parskip} % Para quitar la sangría de los párrafos
\usepackage{fancyhdr}
\usepackage[spanish]{babel}
\usepackage{csquotes}
\usepackage{tikz-cd}

\usepackage{eso-pic} % Permite colocar elementos absolutos en la página

\pagestyle{fancy}

\setlength{\headsep}{7pt}

% Definir un contador que se reinicie con cada sección
\newcounter{defcounter}[section]
\renewcommand{\thedefcounter}{\thesection.\arabic{defcounter}}

\newcommand{\mydef}{
  \refstepcounter{defcounter}
  \textbf{Def. \thedefcounter.}
}

\newcounter{teocounter}[section]
\renewcommand{\theteocounter}{\thesection.\arabic{teocounter}}

\newcommand{\myteo}{
  \refstepcounter{teocounter}
  \textbf{Teo. \theteocounter.}
}

\newcounter{obscounter}[section]
\renewcommand{\theobscounter}{\thesection.\arabic{obscounter}}

\newcommand{\myobs}{
  \refstepcounter{obscounter}
  \textbf{Obs. \theobscounter.}
}

\newcounter{propcounter}[section]
\renewcommand{\thepropcounter}{\thesection.\arabic{propcounter}}

\newcommand{\myprop}{
  \refstepcounter{propcounter}
  \textbf{Prop. \thepropcounter.}
}

\newcounter{lemacounter}[section]
\renewcommand{\thelemacounter}{\thesection.\arabic{lemacounter}}

\newcommand{\mylema}{
  \refstepcounter{lemacounter}
  \textbf{Lema \thelemacounter.}
}

\newcounter{corcounter}[section]
\renewcommand{\thecorcounter}{\thesection.\arabic{corcounter}}

\newcommand{\mycor}{
  \refstepcounter{corcounter}
  \textbf{Cor. \thecorcounter.}
}

\newcounter{excounter}[section]
\renewcommand{\theexcounter}{\thesection.\arabic{excounter}}

\newcommand{\myex}{
  \refstepcounter{excounter}
  \textbf{Ejemplo. \theexcounter.}
}


%--------------------------------------------------------------
\makeatletter
\def\moverlay{\mathpalette\mov@rlay}
\def\mov@rlay#1#2{\leavevmode\vtop{%
   \baselineskip\z@skip \lineskiplimit-\maxdimen
   \ialign{\hfil$\m@th#1##$\hfil\cr#2\crcr}}}
\newcommand{\charfusion}[3][\mathord]{
    #1{\ifx#1\mathop\vphantom{#2}\fi
        \mathpalette\mov@rlay{#2\cr#3}
      }
    \ifx#1\mathop\expandafter\displaylimits\fi}
\makeatother

\newcommand{\cupdot}{\charfusion[\mathbin]{\cup}{\cdot}}
\newcommand{\bigcupdot}{\charfusion[\mathop]{\bigcup}{\cdot}}

\title{\textbf{Resumen MA5802}}
\author{$\,$}
\date{$\,$}

\fancyhf{}
\lhead[\leftmark]{
Facultad de Ciencias Físicas y Matemáticas
}
\rhead{
Universidad de Chile
}

\cfoot{\thepage}

%-----------------------------------------------------------
\begin{document}

\twocolumn[

\textbf{Departamento de Ingeniería Matemática\\
MA5802 - Introducción a los Sistemas Dinámicos y la Teoría Ergódica\\
Profesores:} Alejandro Maass S., Sebastián Donoso F.\\
\textbf{Autor:} Javier Cruz\\


% Colocar imagen en la esquina superior derecha (ajusta las coordenadas si es necesario)
\AddToShipoutPictureBG*{%
  \put(\LenToUnit{\paperwidth - 5.25cm},\LenToUnit{\paperheight - 3.1cm}){%
    \includegraphics[width=3.25cm]{img/fcfm2.pdf}%
  }%
}

\begin{center}
    \textbf{\LARGE Resumen Control 1}
\end{center}
\vspace{0.6cm}
]

\section{Recuerdo}
\hrule

\myteo (Carathéodory):\\
Sea \(\mathcal{S}\subseteq \mathcal{P}(X)\) una semi\(-\)álgebra, y 
\(\mu: \mathcal{S} \rightarrow \overline{\mathbb{R}}_+\) una medida. 

Si definimos \(\mu^*: \mathcal{P}(X)\rightarrow \overline{\mathbb{R}}_+\) por

\[
\mu^*(A) = \inf_{(A_n) \in \mathcal{R}(A)} \sum_n \mu(A_n),
\]

entonces
\begin{enumerate}[label=\roman*.]
    \item \(\mu^*\) es una medida exterior.
    \item \(\mu^*\) es una extensión de \(\mu\).
    \item La \(\sigma-\)álgebra \(\eta\) de conjuntos \(\mu^*-\)medibles contiene a
    \(\sigma(\mathcal{S})\).
\end{enumerate}


\shorthandoff{"}

\section{Introducción}
\hrule

\mydef (Sistema Dinámico Abstracto):\\
Sea \((X,\mathcal{B},\mu)\) un espacio de medida de probabilidad, y 
\(T: X \rightarrow X\) una función \(\mathcal{B}-\mathcal{B}\) medible, 
\textit{biyectiva} y \(T\mu = \mu\) (i.e., \textit{preserva la medida}), donde

\[
T\mu (B) = \mu (T^{-1}(B)) = \mu (B), \quad B \in \mathcal{B}.
\]

Se dice que \((X,\mathcal{B},\mu,T)\) es un \textit{sistema dinámico abstracto} (s.d.a.) 
o un \textit{sistema dinámico que preserva la medida} \(\mu\) (s.p.m.).

\mydef (Factor o Extensión):\\
Se dice que el s.d.a. \((Y,\mathcal{A},\nu,S)\) es \textit{factor} del s.d.a. 
\((X,\mathcal{B},\mu,T)\) (o que \((X,\mathcal{B},\mu,T)\) es una \textit{extensión} de 
\((Y,\mathcal{A},\nu,S)\)) ssi \(\exists \psi : X \rightarrow Y\), 
\(\mathcal{B}-\mathcal{A}-\)medible, sobreyectiva tal que \(\psi \mu = \nu\) (i.e., 
\(\psi\mu(A) = \mu(\psi^{-1}(A))=\nu(A)\), con \(A \in \mathcal{A}\)) y 
\(S \circ \psi = \psi \circ T\).

% Figura que se puede borrar
\[
\begin{tikzcd}
X \arrow[r, "T"] \arrow[d, "\psi"] & X \arrow[d, "\psi"] \\
Y \arrow[r, "S"]                   & Y
\end{tikzcd}
\]

\mydef (Conjugación o Isomorfía):\\
Se dice que los s.d.a. son \textit{isomorfos} ssi \(\psi\) es biyectiva y 
\textit{bi}-medible (i.e., \(\psi\) y \(\psi^{-1}\) son medibles).

\myex (Rotación Irracional):\\
Sean
\begin{itemize}
    \item \(X = \mathbb{S}^1 = \{z \in \mathbb{C}: |z|=1\}\).
    \item \(\alpha \in \mathbb{R}\setminus\mathbb{Q}\).
    \item \(T:X \rightarrow X,\,z \mapsto T(z) = e^{i2\pi\alpha}\cdot z\).
\end{itemize}

Si \(\mathcal{B}(X)\) es la \(\sigma-\)álgebra de los borelianos de \(X\) y \(\mu\)
es la medida uniforme en \(\mathbb{S}^1\) (o de Haar, o de Lebesgue) se tiene que 
\((X,\mathcal{B},\mu,T)\) es un s.d.a.

\section{Recurrencia}
\hrule

\myteo (Poincaré Débil):\\
Si \((X,\mathcal{B},\mu,T)\) es s.d.a. y \(B \in \mathcal{B}\) con \(\mu(B)>0\),
entonces existe \(n>1\) tal que 

\[
\mu(B \cap T^{-n}B) = \mu(\{x \in B: T^nB \in B\})>0.
\]

\input{classes/2}

\input{classes/3}

\input{classes/4}

\input{classes/5}

\input{classes/6}

\input{classes/7}

\end{document}